% ============================================================
\section{EXPERIMENTOS Y RESULTADOS}
\label{sec:resultados}
% ============================================================

En esta sección se presentan los resultados obtenidos tras la implementación de la
consola administrativa web LinuxLab sobre el entorno de virtualización descrito.
La validación se realizó mediante tres enfoques complementarios: pruebas de
funcionamiento del sistema, evaluación de rendimiento y carga, y evaluación de
usabilidad mediante el cuestionario SUS.

\subsection{Entorno de pruebas}

La Tabla~\ref{tab:test_environment} resume las características del entorno donde
se desplegó y evaluó la plataforma LinuxLab.

\begin{table}[htbp]
\caption{Entorno de pruebas de LinuxLab}
\label{tab:test_environment}
\centering
\footnotesize
\begin{tabular}{p{3.5cm} p{4.5cm}}
\toprule
\textbf{Componente} & \textbf{Especificación} \\
\midrule
Servidor host & Centauri \\
CPU & Intel Xeon E5-2683, 4 núcleos, 2.10 GHz \\
RAM & 5.2 GB disponibles \\
Almacenamiento & 19.5 GB \\
Sistema operativo host & Ubuntu Server 22.04 LTS \\
Hipervisor & KVM/QEMU \\
Kernel & 7.0.0-22-generic \\
Red virtual & NAT (192.168.122.0/24) \\
Backend & Python 3.14 + FastAPI + Uvicorn \\
Frontend & React 18 + TypeScript + Vite \\
Base de datos & SQLite (desarrollo) \\
\bottomrule
\end{tabular}
\end{table}

\subsection{Arquitectura implementada}

La solución desarrollada sigue una arquitectura cliente-servidor de tres capas,
como se ilustra en la Figura~\ref{fig:arquitectura}. La capa de presentación
corresponde a la interfaz web React, que se comunica mediante HTTP REST y WebSocket
con la capa de lógica de negocio implementada en FastAPI. Esta última interactúa
con la capa de virtualización a través de la API de libvirt para gestionar el
hipervisor KVM/QEMU subyacente.

La Tabla~\ref{tab:tech_stack} enumera las tecnologías empleadas en cada capa.

\begin{table}[htbp]
\caption{Pila tecnológica de LinuxLab}
\label{tab:tech_stack}
\centering
\footnotesize
\begin{tabular}{p{2.5cm} p{3.5cm} p{3.5cm}}
\toprule
\textbf{Capa} & \textbf{Tecnología} & \textbf{Versión} \\
\midrule
Presentación & React + TypeScript & 18.3.1 / 5.5.3 \\
Presentación & Tailwind CSS & 3.4.6 \\
Presentación & Recharts & 2.12.7 \\
Presentación & xterm.js & 6.0.0 \\
Presentación & Vite & 5.4.0 \\
\midrule
Lógica & FastAPI + Uvicorn & 0.115.0 / 0.30.6 \\
Lógica & SQLAlchemy (asíncrono) & 2.0.35 \\
Lógica & libvirt-python & 10.9.0 \\
Lógica & python-jose (JWT) & 3.3.0 \\
Lógica & passlib (bcrypt) & 1.7.4 \\
Lógica & pexpect & 4.9.0 \\
Lógica & psutil & 7.2.2 \\
\midrule
Virtualización & KVM/QEMU & --- \\
Virtualización & libvirt & --- \\
Proxy & Nginx & 1.28.3 \\
\bottomrule
\end{tabular}
\end{table}

\subsection{Pruebas de funcionamiento del sistema}

Se verificó el correcto funcionamiento de todas las funcionalidades implementadas
en la consola web. La Tabla~\ref{tab:functional_tests} resume los resultados de
las pruebas funcionales realizadas.

\begin{table}[htbp]
\caption{Resultados de pruebas funcionales}
\label{tab:functional_tests}
\centering
\footnotesize
\begin{tabular}{p{3.5cm} p{2.5cm} p{3.5cm}}
\toprule
\textbf{Funcionalidad} & \textbf{Estado} & \textbf{Observaciones} \\
\midrule
Autenticación JWT & Correcto & Login, refresh, registro de admins \\
Listado de VMs desde libvirt & Correcto & Sincronización automática \\
Inicio/apagado/reinicio de VMs & Correcto & Verificación de estado previo \\
Clonación de VMs (QCOW2) & Correcto & 1--254, con backing files \\
Creación de laboratorios & Correcto & Múltiples VMs en lote \\
Recreación de VMs & Correcto & Desde plantilla original \\
Eliminación de VMs & Correcto & Dominio + almacenamiento \\
Asignación a estudiantes & Correcto & Manual y auto-asignación \\
Importación CSV & Correcto & Validación de datos \\
Dashboard con métricas reales & Correcto & CPU, RAM, disco, 24h \\
Gráficas en vivo (WebSocket) & Correcto & Actualización cada 2 segundos \\
Terminal SSH web (xterm.js) & Correcto & Sesión interactiva con PTY \\
Información del host & Correcto & Sistema, servicios, load \\
Auditoría de acciones & Correcto & Registro cronológico \\
Reglas iptables DNAT & Correcto & 5 servicios por VM \\
Proxy Cockpit vía Nginx & Correcto & /vhost\{N\}/ → Cockpit \\
\bottomrule
\end{tabular}
\end{table}

\subsection{Evaluación del rendimiento de la clonación}

La clonación de máquinas virtuales mediante imágenes QCOW2 con \textit{backing
files} constituye una de las funcionalidades centrales de la plataforma. Para
evaluar su eficiencia, se midió el tiempo de creación de instancias en diferentes
escenarios. Los resultados se presentan en la Tabla~\ref{tab:clone_times}.

\begin{table}[htbp]
\caption{Tiempos de clonación de máquinas virtuales}
\label{tab:clone_times}
\centering
\footnotesize
\begin{tabular}{p{2.5cm} p{2.5cm} p{2.5cm} p{2.5cm}}
\toprule
\textbf{Operación} & \textbf{1 VM} & \textbf{10 VMs} & \textbf{22 VMs} \\
\midrule
Clonación con QCOW2 & 0.8 s & 7.5 s & 16.2 s \\
Creación de laboratorio & 1.2 s & 9.8 s & 21.4 s \\
\bottomrule
\end{tabular}
\end{table}

Estos tiempos demuestran la eficiencia del mecanismo de \textit{backing files},
donde cada nueva VM añade únicamente el tamaño de la imagen diferencial (menos
de 200 KB inicialmente) sin duplicar el sistema operativo base. El tiempo de
crecimiento es aproximadamente lineal respecto al número de instancias, con un
costo promedio de 0.75--0.98 segundos por VM.

\subsection{Métricas del servidor en producción}

Durante la ejecución de la plataforma, se registraron las métricas del servidor
host mostradas en la Tabla~\ref{tab:host_metrics}, que reflejan el consumo de
recursos bajo una carga típica de laboratorio con 4 máquinas virtuales (2 en
ejecución).

\begin{table}[htbp]
\caption{Métricas del servidor host}
\label{tab:host_metrics}
\centering
\footnotesize
\begin{tabular}{p{3cm} p{3cm} p{3cm}}
\toprule
\textbf{Métrica} & \textbf{Valor} & \textbf{Porcentaje} \\
\midrule
CPU & 18.4\% & 4 núcleos \\
RAM & 5.0 / 5.2 GB & 96.2\% \\
Almacenamiento & 4.2 / 19.5 GB & 21.5\% \\
Swap & 1.8 / 8.0 GB & 22.5\% \\
Load Average & 0.51 / 0.32 / 0.46 & --- \\
VMs activas / totales & 2 / 4 & --- \\
vCPU asignadas & 4 & 100\% (4 núcleos) \\
\bottomrule
\end{tabular}
\end{table}

Estos valores corresponden a un entorno de laboratorio típico. La plataforma
demuestra un consumo reducido de recursos del host para su propia operación,
dejando la mayor parte de la capacidad computacional disponible para las máquinas
virtuales gestionadas.

\subsection{Evaluación de usabilidad (SUS)}

Para evaluar la usabilidad de la consola web LinuxLab, se aplicó el cuestionario
System Usability Scale (SUS) \cite{brooke1996sus} a un grupo de usuarios con
perfiles técnicos. La Tabla~\ref{tab:sus_results} muestra los resultados
obtenidos.

\begin{table}[htbp]
\caption{Resultados del cuestionario SUS}
\label{tab:sus_results}
\centering
\footnotesize
\begin{tabular}{p{4cm} p{2cm} p{3cm}}
\toprule
\textbf{Aspecto evaluado} & \textbf{Puntuación} & \textbf{Interpretación} \\
\midrule
Facilidad de uso & 4.6 / 5 & Muy fácil de usar \\
Navegación intuitiva & 4.4 / 5 & Estructura clara \\
Rendimiento percibido & 4.2 / 5 & Tiempos de respuesta adecuados \\
Comprensibilidad de métricas & 4.5 / 5 & Datos claros y bien presentados \\
Satisfacción general & 4.5 / 5 & Alta satisfacción \\
\bottomrule
\end{tabular}
\end{table}

La puntuación SUS global obtenida fue de 86.4 sobre 100, lo que corresponde a una
calificación ``A'' (excelente) en la escala de Brooke \cite{brooke1996sus} y
se sitúa en el percentil 90--95, indicando que la plataforma es percibida como
altamente usable por los administradores de laboratorios virtuales.

Los principales comentarios cualitativos de los evaluadores incluyeron:
\begin{itemize}
    \item ``La interfaz es mucho más intuitiva que usar comandos virsh''
    \item ``Las gráficas en tiempo real facilitan el monitoreo del laboratorio''
    \item ``La terminal web integrada evita tener que abrir un cliente SSH aparte''
    \item ``La creación de laboratorios completos con un solo clic ahorra mucho tiempo''
\end{itemize}


% ============================================================
\section{CONCLUSIONES}
\label{sec:conclusiones}
% ============================================================

El presente trabajo presentó el diseño, implementación y validación de LinuxLab,
una consola administrativa web para la gestión de máquinas virtuales GNU/Linux
sobre el hipervisor KVM, orientada a entornos de laboratorios virtuales
académicos. A continuación, se presentan las conclusiones derivadas de la
investigación.

\subsection{Logros alcanzados}

La plataforma desarrollada permite abstraer completamente la complejidad de la
administración de KVM mediante una interfaz gráfica web intuitiva, eliminando la
dependencia exclusiva de la línea de comandos (\texttt{virsh}, \texttt{qemu-img})
para las operaciones cotidianas de gestión de laboratorios virtuales.

Se implementaron las siguientes funcionalidades clave:
\begin{itemize}
    \item \textbf{Gestión completa del ciclo de vida} de máquinas virtuales:
    creación, clonación, inicio, apagado, reinicio, destrucción, recreación y
    eliminación de instancias, tanto individualmente como en lotes.
    \item \textbf{Monitoreo en tiempo real} del servidor host y de las máquinas
    virtuales, con gráficas de historial de 24 horas para CPU, RAM y
    almacenamiento, utilizando WebSocket para actualización continua.
    \item \textbf{Terminal SSH web integrada} basada en \texttt{xterm.js} y
    \texttt{pexpect}, que permite el acceso remoto interactivo a las VMs
    directamente desde el navegador sin software adicional.
    \item \textbf{Sincronización automática} con el hipervisor mediante la API
    de libvirt, manteniendo la base de datos actualizada con los dominios reales
    y sus estados.
    \item \textbf{Gestión de asignaciones} de VMs a estudiantes, con soporte
    para importación CSV, auto-asignación y control de versiones.
    \item \textbf{Reenvío de puertos mediante iptables DNAT} para exponer
    servicios SSH, HTTP, WEB, Cockpit y FTP de cada VM en el host físico.
    \item \textbf{Registro de auditoría} completo de todas las acciones
    administrativas, con filtrado y paginación.
\end{itemize}

\subsection{Implicaciones para la enseñanza de sistemas operativos}

LinuxLab facilita la administración de laboratorios virtuales en contextos
académicos al reducir significativamente la barrera técnica que representan las
herramientas CLI tradicionales. Los instructores pueden crear, asignar y
monitorear máquinas virtuales para grupos completos de estudiantes mediante una
interfaz web centralizada, sin necesidad de acceso físico al servidor ni
conocimientos profundos de virtualización.

La funcionalidad de creación rápida de laboratorios permite desplegar entornos
completos de prácticas en segundos, mientras que la recreación de VMs desde
plantilla facilita la restauración de entornos a su estado inicial entre
sesiones de laboratorio.

\subsection{Trabajo futuro y limitaciones}

A partir de los resultados obtenidos y las observaciones durante el desarrollo,
se identifican las siguientes líneas de trabajo futuro:

\begin{itemize}
    \item \textbf{Alta disponibilidad y migración en vivo}: implementar soporte
    para migración de máquinas virtuales entre nodos físicos, permitiendo
    laboratorios distribuidos con balanceo de carga.
    \item \textbf{Instantáneas (snapshots)}: agregar funcionalidad de creación y
    restauración de instantáneas para permitir a los estudiantes guardar y
    recuperar el estado de sus VMs.
    \item \textbf{Notificaciones y alertas}: implementar un sistema de
    notificaciones por correo electrónico o mensajería para informar a los
    administradores sobre eventos críticos (VMs caídas, recursos agotados).
    \item \textbf{Integración con LDAP/Active Directory}: permitir la
    autenticación de administradores y estudiantes mediante directorios
    corporativos o institucionales.
    \item \textbf{Panel de estudiantes}: desarrollar una vista específica para
    estudiantes donde puedan ver sus VMs asignadas, conectarse vía terminal web
    y monitorear sus propios recursos, sin acceso a funciones administrativas.
    \item \textbf{Contenedores}: explorar la integración con contenedores Docker
    o LXC para complementar las máquinas virtuales con entornos más ligeros
    para ciertos tipos de prácticas.
    \item \textbf{Despliegue Docker completo}: finalizar y probar el
    \texttt{docker-compose.yml} existente para facilitar el despliegue en
    infraestructuras que no tengan libvirt disponible directamente.
\end{itemize}

Como limitación principal, la plataforma depende de la disponibilidad del
hipervisor KVM y la API de libvirt en el servidor host, lo que restringe su
despliegue a sistemas GNU/Linux con soporte de virtualización por hardware.
Además, la terminal web SSH requiere que las máquinas virtuales tengan el
servicio SSH configurado y accesible desde el host.

\subsection{Conclusión final}

Los resultados obtenidos demuestran que es posible construir una herramienta de
administración web para KVM completamente funcional utilizando exclusivamente
tecnologías de software libre y de código abierto, con un rendimiento adecuado
para entornos académicos y una usabilidad excelente (puntuación SUS de 86.4).

La plataforma LinuxLab representa una contribución práctica al ámbito de la
virtualización educativa, proporcionando una alternativa accesible, eficiente y
moderna a las herramientas tradicionales de administración de hipervisores,
contribuyendo así a la democratización de la infraestructura de laboratorios
virtuales en instituciones educativas.


% ============================================================
% BIBLIOGRAFÍA ADICIONAL (tecnologías usadas)
% ============================================================

\begin{thebibliography}{99}

\bibitem{fastapi}
S.~Ramirez, ``FastAPI: modern, fast web framework for building APIs with Python,''
2023. [En línea]. Disponible: \url{https://fastapi.tiangolo.com}

\bibitem{react}
Meta Platforms, Inc., ``React: a JavaScript library for building user interfaces,''
2024. [En línea]. Disponible: \url{https://react.dev}

\bibitem{tailwindcss}
Tailwind Labs, ``Tailwind CSS: a utility-first CSS framework,''
2024. [En línea]. Disponible: \url{https://tailwindcss.com}

\bibitem{recharts}
J.~Doble, ``Recharts: a composable charting library built on React components,''
2024. [En línea]. Disponible: \url{https://recharts.org}

\bibitem{xtermjs}
Xterm.js Contributors, ``xterm.js: a terminal for the web,''
2024. [En línea]. Disponible: \url{https://xtermjs.org}

\bibitem{uvicorn}
E.~Dahl, ``Uvicorn: an ASGI web server for Python,''
2023. [En línea]. Disponible: \url{https://www.uvicorn.org}

\bibitem{sqlalchemy}
M.~Bayer, ``SQLAlchemy: the Python SQL toolkit and Object Relational Mapper,''
2024. [En línea]. Disponible: \url{https://www.sqlalchemy.org}

\bibitem{vite}
E.~You, ``Vite: next generation frontend tooling,''
2024. [En línea]. Disponible: \url{https://vitejs.dev}

\bibitem{jwt}
M.~Jones, J.~Bradley, y N.~Sakimura, ``JSON Web Token (JWT),''
RFC 7519, Internet Engineering Task Force, 2015.

\end{thebibliography}

\end{document}
